\begin{abstract}
Web agents act through chains of tool calls over REST and Model Context Protocol (MCP) services. In a ReAct loop, every dependent call costs an LLM round over a growing context. Letting the agent write code removes these rounds only for mechanical steps, such as looking up a field or computing a value. Semantic steps still send control back to the model: deciding which returned entity the user meant, which branch applies, or which items to keep. We present \jif, a \emph{graph-guided tool compiler}. The agent writes the whole tool chain once as a program, and the chain then runs to completion without returning to the LLM.

The compiler rests on two designs. First, it reads a \emph{split object--operation graph} (\sog). The object graph describes business entities, properties, relations and rules. The operation graph describes tool-backed operations, and its data-flow edges are derived from which operations produce and consume each property. From the \sog, the compiler renders typed stubs for the agent, statically checks the agent's program, and obtains the entity semantics that run-time decisions need. Second, every semantic decision in the program is written as an explicit conditional on a System-1 primitive (\texttt{jev.choice}, \texttt{jev.noul}). These conditionals are evaluated at run time by a calibrated non-generative model (Jev), and a confidence gate hands a decision back to the LLM only when it is uncertain.

The agent therefore writes code once, the compiler checks and optimizes it, and Jev evaluates its decisions. The LLM is called to plan, to answer, and again only when a decision is uncertain. \todo{Headline results on FinMCP-Bench vs.\ ReAct / ReWOO / LLMCompiler / programmatic tool calling: LLM calls, tokens, latency at no loss of success; reduction of invalid-call and wrong-entity errors; main effects of graph guidance and \jif.}
\end{abstract}
